\documentclass[11pt]{article}

\usepackage[margin=0.9in]{geometry}
\usepackage{amsmath,amssymb}
\usepackage{booktabs}
\usepackage{longtable}
\usepackage{array}
\usepackage{enumitem}
\usepackage[T1]{fontenc}
\usepackage[utf8]{inputenc}
\usepackage[hidelinks]{hyperref}

% --- convenience macros -------------------------------------------------
\newcommand{\High}{\textbf{High}}
\newcommand{\Med}{Medium}
\newcommand{\Low}{\textbf{Low}}
\newcommand{\keywords}[1]{\par\smallskip\noindent\textit{Keywords:} #1}
\newcommand{\rating}[1]{\par\smallskip\noindent\textit{Rating justification (mine):} #1}

\title{Fluid Metrics Exploration Tracker\\[4pt]
}
\author{Hooligans}
\date{August 2026}

\begin{document}
\maketitle

\medskip
\noindent
The meteorological literature contains an
explicit statement of our shock-shift problem under the name \emph{double penalty}
\cite{wernli2008}; the relationship between the negative Sobolev norm and the Wasserstein
distance is a quantitative two-sided bound with stated conditions, not a vague analogy
\cite{peyre2018}; the ``wake number'' item has a direct, practical implementation in the
visualization literature under the name \emph{winding angle} \cite{sadarjoen2000}; and
Arnold's paper does in fact prove that the sectional curvature of the volume-preserving
diffeomorphism group is negative in most sections \cite{arnold1966}. Two ratings changed as
a result, and one item remains genuinely unsupported.

\tableofcontents

%=========================================================================
\section{How to use this document}

\begin{enumerate}[leftmargin=*]
\item The \textbf{master table} in Section~\ref{sec:master} is the working checklist.
Update the \texttt{Status} column as items move from \texttt{Not tested} to
\texttt{In progress} to \texttt{Tested (result)}, and fill in \texttt{Owner}.
\item The theme sections give, for each item: what it actually is, where to look, search
keywords, whether the blocker is a mathematical gap or an implementation gap, known
pitfalls, and a promise rating with justification.
\item \textbf{All promise ratings are my judgment, not literature consensus.} They are
explained so that you can disagree with the reasoning rather than with the label.
\end{enumerate}

\subsection{Rating rubric (my definitions)}

\begin{longtable}{@{}p{0.13\textwidth}p{0.82\textwidth}@{}}
\toprule
\textbf{Rating} & \textbf{Meaning} \\
\midrule
\endhead
\High{} & Strong precedent in an adjacent field solving the \emph{same} structural
problem, mature software exists, and the metric targets a failure mode that $L^2$ provably
misses. Test in the first wave. \\
\addlinespace
\Med{} & Plausible and interesting, but either the precedent is indirect, the
implementation is nontrivial, or the metric likely duplicates information that another
metric already provides. Test after the first wave. \\
\addlinespace
\Low{} & Beautiful or deep, but the theory is immature, the tooling is research grade, or
the quantity is too indirect to serve as a routine diagnostic within the next year. Park
it, and revisit if a first-wave metric fails in a way that this item would explain. \\
\bottomrule
\end{longtable}


%=========================================================================
\section{Master tracker table}
\label{sec:master}

Column meanings: \emph{Gap} records whether the blocker is mathematical (M),
implementation (I), or both; \emph{Cost} is a rough estimate of effort and compute;
\emph{Promise} is my rating per the rubric above. Two ratings changed in this revision and
are marked with a dagger.

\begin{longtable}{@{}p{0.055\textwidth}p{0.40\textwidth}p{0.045\textwidth}p{0.045\textwidth}p{0.09\textwidth}p{0.13\textwidth}p{0.08\textwidth}@{}}
\toprule
\textbf{ID} & \textbf{Item} & \textbf{Gap} & \textbf{Cost} & \textbf{Promise} &
\textbf{Status} & \textbf{Owner} \\
\midrule
\endfirsthead
\toprule
\textbf{ID} & \textbf{Item} & \textbf{Gap} & \textbf{Cost} & \textbf{Promise} &
\textbf{Status} & \textbf{Owner} \\
\midrule
\endhead
\multicolumn{7}{@{}l}{\emph{Optimal transport and displacement-aware distances}}\\
OT-1 & Wasserstein ($W_1$, $W_2$) for horizontal displacement & M & Med & \High{} & Not tested & \\
OT-2 & Signed-field fix: graph-space OT, unbalanced OT, $\pm$ split & M & Med & \High{} & Not tested & \\
OT-3 & Sliced Wasserstein & I & Low & \High{} & Not tested & \\
OT-4 & Entropic (Sinkhorn) relaxation and Kantorovich dual & I & Low & \Med{} & Not tested & \\
OT-5 & $W_1$ on increment and gradient PDFs, not on fields & I & Low & \High{} & Not tested & \\
\addlinespace
\multicolumn{7}{@{}l}{\emph{Shock geometry and detection}}\\
SH-1 & Derivative and sharpness shock detector & I & Low & \High{} & Not tested & \\
SH-2 & Distance between extracted shock sets (ASSD, Hausdorff-95) & I & Low & \High{} & Not tested & \\
SH-3 & Shock mask overlap (IoU, Dice, boundary F-score) & I & Low & \Med{} & Not tested & \\
SH-4 & Shock front position and speed error; Rankine--Hugoniot residual & M+I & Med & \Med{} & Not tested & \\
\addlinespace
\multicolumn{7}{@{}l}{\emph{Norms and function-space metrics}}\\
NM-0 & Pointwise $L^p$ baselines (MAE, MSE, RMSE, NRMSE) & --- & Low & n/a$^\ddagger$ & Implemented; measured & \\
NM-1 & Sobolev $H^s$ norm with $s>0$, that is $H^1$ & --- & Low & \Low{} & Not tested & \\
NM-2 & Negative Sobolev and dual norms, $\dot H^{-1}$ and $H^{-s}$ & M & Low & \High{} & Not tested & \\
NM-3 & Mollification split $f = f*\eta_\varepsilon + \Delta$, metricize parts separately & M & Low & \High{} & Not tested & \\
\addlinespace
\multicolumn{7}{@{}l}{\emph{Curvature and differential geometry}}\\
CG-1 & Level-set and isosurface curvature statistics & I & Low & \Med{} & Not tested & \\
CG-2 & Edge detection (Canny, Sobel) as a cheap shock localizer & I & Low & \Med{} & Not tested & \\
CG-3 & Discrete differential geometry and DEC on meshes & M & Med & \Low{} & Not tested & \\
CG-4 & Metric-geometry curvature (Ollivier--Ricci, Alexandrov) & M & Med & \Low{} & Not tested & \\
CG-5 & Riemann curvature (Arnold geometric hydrodynamics) & M & High & \Low{} & Not tested & \\
\addlinespace
\multicolumn{7}{@{}l}{\emph{Vorticity and physics invariants}}\\
PH-1 & Entropy production, signed and three-way & M+I & Low & \High{} & Not tested & \\
PH-2 & PDE residual in weak, finite-volume form & I & Low & \High{} & Not tested & \\
PH-3 & Solver-consistency residual (one step of a trusted solver) & I & Med & \High{} & Not tested & \\
PH-4 & Enstrophy & I & Low & \Med{} & Implemented; measured$^\S$ & \\
PH-5 & Palinstrophy & I & Low & \Med{} & Not tested & \\
PH-6 & Helmholtz split: ``take curls but protect from shocks'' & M & Med & \High{} & Not tested & \\
\addlinespace
\multicolumn{7}{@{}l}{\emph{Pattern and feature detection}}\\
PD-1 & Hough transform for eddy size & I & Med & \Low{}$^\dagger$ & Not tested & \\
PD-2 & Winding angle and Poincar\'e index (``wake number'') & I & Med & \High{}$^\dagger$ & Not tested & \\
PD-3 & Eddy count distribution and eddy census & I & Med & \Med{} & Not tested & \\
PD-4 & $Q$-criterion, $\lambda_2$, Okubo--Weiss as eddy detectors & I & Low & \Med{} & Not tested & \\
\addlinespace
\multicolumn{7}{@{}l}{\emph{Topological data analysis}}\\
TD-1 & One-parameter persistence diagrams (sublevel sets) & I & Low & \High{} & Not tested & \\
TD-2 & Two-parameter (multiparameter) persistence & M & High & \Low{} & Not tested & \\
TD-3 & Well groups & M & High & \Low{} & Not tested & \\
TD-4 & Reeb graphs and Morse--Smale complexes & I & Med & \Med{} & Not tested & \\
TD-5 & Jacobi sets (native two-field comparison) & I & Med & \Med{} & Not tested & \\
\addlinespace
\multicolumn{7}{@{}l}{\emph{Basis decompositions}}\\
BD-1 & Fourier transform and energy spectrum & --- & Low & \Med{} & Not tested & \\
BD-2 & Two-point correlation & --- & Low & \Med{} & Not tested & \\
BD-3 & Wavelets and coherent vortex extraction (CVS) & I & Med & \High{} & Not tested & \\
BD-4 & Wavelet scattering and phase harmonics (``turbulence transform'') & I & Med & \High{} & Not tested & \\
\addlinespace
\multicolumn{7}{@{}l}{\emph{Probabilistic and distributional metrics}}\\
PS-1 & Energy distance & I & Low & \Med{} & Not tested & \\
PS-2 & CRPS & I & Low & \High{}$^{*}$ & Not tested & \\
PS-3 & Spread-to-skill ratio & I & Low & \High{}$^{*}$ & Not tested & \\
PS-4 & Heavy-tail quantifiers (flatness, Hill estimator, EVT) & I & Low & \High{} & Not tested & \\
PS-5 & PDF of pairwise distances, Ripley's $K$ & I & Med & \Med{} & Not tested & \\
\addlinespace
\multicolumn{7}{@{}l}{\emph{Infrastructure and protocol}}\\
IN-1 & Project onto planes and compare & I & Low & \Med{} & Not tested & \\
IN-2 & Adaptive meshes, comparison across meshes & I & Med & \Med{} & Implemented (remap) & \\
IN-3 & Degradation-ladder protocol with Spearman acceptance test & I & Med & \High{} & Implemented & \\
IN-4 & Gaussian-impostor canary & I & Low & \High{} & Implemented; measured & \\
\bottomrule
\end{longtable}

\noindent
$^{*}$ High \emph{only if} ensembles are available. For a deterministic surrogate these
two items are not applicable and should be marked N/A rather than left as untested.

\noindent
$^\ddagger$ NM-0 is not a candidate and carries no promise rating. It is the pointwise
control that every other entry is read against; it is implemented so that each candidate has
a baseline to be compared with. See Section~\ref{sec:nm0}.

\noindent
$^\S$ PH-4 is implemented and measured, and the measurement disqualifies it as a
\emph{comparison} metric: a translated field has identical enstrophy by construction, so its
clean and unrelated-field values coincide and it has no dynamic range on the normalised
scale. It remains usable as a single-field tripwire, which is how this document already
describes it.

\noindent
$^\dagger$ Rating changed in this revision. PD-2 rose from Medium to High because the
winding angle turns out to be a standard, implemented vortex-detection method with
documented strengths and weaknesses \cite{sadarjoen2000}. PD-1 fell from Medium to Low
because the eddy-detection precedent I expected to find for the Hough transform does not
appear to exist, and the winding-angle method covers the same ground with actual support.

%=========================================================================
\section{Optimal transport and displacement-aware distances}

This is the theme where the meeting's central insight, that horizontal distances make
sense while vertical distances do not, has the strongest outside precedent. Two separate
communities hit exactly this wall and solved it, and one of them gave the pathology a name.

\subsection{OT-1. Wasserstein for horizontal displacement --- \High{}}

\paragraph{What it is.} Instead of asking how different two fields are at each point, one
asks what the cheapest way is to move mass from one field's profile to the other's. A pure
translation of a shock by a distance $s$ then costs an amount that grows with $s$, rather
than saturating the way $L^2$ does.

\paragraph{Why the precedent matters.} Seismic full-waveform inversion suffers an identical
pathology, known there as cycle skipping: the $L^2$ misfit between a predicted and an
observed waveform behaves badly when the predicted arrival is shifted in time. Engquist and
Froese proposed the Wasserstein metric as the remedy, observing that seismic signals are
conventionally compared either by travel-time difference or by $L^2$ difference and that
the Wasserstein metric ``exhibits properties from both of the traditional measures''
\cite{engquist2014}. The convexity property is precisely the one we want: in the review
literature the $W_2$ misfit ``precisely captures the data shift $s$ as the misfit is $s^2$
independent of the profile of the signal'' \cite{engquistyang}. That is a clean statement
that $W_2$ measures \emph{position} error while remaining blind to \emph{profile shape},
which is exactly the decomposition the meeting was reaching for.

\paragraph{The math gap you identified is a known, named problem.} Converting a pressure or
velocity field into a probability distribution requires nonnegativity and a fixed total
mass. Yang, Engquist, Sun and Froese devote explicit attention to ``the difficulty of
transforming seismic signals into nonnegative density functions'' \cite{yang2016}, and
M\'etivier and colleagues state the obstruction plainly: ``the main difficulty for the
application of optimal transport to seismic data is related to their non positivity: the
optimal transport theory is developed for the comparison of positive quantities''
\cite{metivier2018eage}. The standard fixes are catalogued under OT-2.

\keywords{quadratic Wasserstein misfit, cycle skipping, optimal transport waveform
inversion, Monge--Amp\`ere, displacement convexity, $W_2$ convexity with respect to
translation.}

\paragraph{Pitfalls.} On a normalized field, $W_2$ is sensitive only to where mass sits, so
it will accept a field with the wrong amplitude but the right position. It must therefore be
paired with an amplitude-sensitive metric. In two and three dimensions the ``horizontal
shift'' interpretation is also less clean than it is in one dimension.

\rating{High, because the structural analogy is exact, the seismic community has a decade of
experience with the failure modes, and the metric directly attacks the shock-position
problem that $L^2$ cannot see.}

\subsection{OT-2. Signed-field fix --- \High{}}

\paragraph{What it is.} There are four standard escapes from the positivity and
mass-normalization constraints. First, split the field into positive and negative parts and
transport each separately. Second, add a constant and renormalize, which is crude and
introduces a tunable parameter. Third, use \emph{unbalanced} optimal transport, in the
Hellinger--Kantorovich sense, which permits mass creation and destruction at a price.
Fourth, use \emph{graph-space} optimal transport, which compares a discrete graph of the
data rather than the data themselves, and thereby preserves convexity with respect to
shifts while handling signed, non-normalized data natively \cite{metivier2018geophysics}.
The graph-space approach grew out of the same group's earlier optimal-transport misfit for
seismograms \cite{metivier2016gji}, and has since been generalized as a family of distances
that contains the $L^p$ distances as special cases \cite{metivier2019ip}.

\rating{High, because this is the single item that unblocks OT-1 for real pressure and
velocity fields. It is a mathematical gap in your taxonomy, and it is a gap that the
seismic community has already filled, so we would be importing rather than inventing.}

\keywords{unbalanced optimal transport, Hellinger--Kantorovich, graph-space OT, partial
optimal transport, signed measure transport.}

\subsection{OT-3. Sliced Wasserstein --- \High{}}

\paragraph{What it is.} Project the field onto many random one-dimensional directions,
compute the cheap sort-based one-dimensional Wasserstein distance along each, and average.
The cost is close to linear.

\paragraph{Note the convergence on your own boards.} ``Project onto planes and compare''
from the Winging It board, catalogued here as IN-1, \emph{is} sliced Wasserstein whenever
the comparison on each projection is a one-dimensional optimal-transport distance. Two
independently generated ideas are the same idea, and they should be treated as a single
work item.

\keywords{sliced Wasserstein distance, max-sliced Wasserstein, Radon transform,
projection-based optimal transport.}

\paragraph{Pitfalls.} Slicing discards information in a way that is difficult to
characterize, so the meeting's remark that it ``introduces its own errors'' is correct. Use
it for screening and for cost-scaling studies, and validate it against exact Wasserstein
distances on small cases.

\subsection{OT-4. Entropic relaxation and Sinkhorn --- \Med{}}

This becomes relevant only \emph{after} OT-1 has demonstrated utility, which matches the
meeting's own sequencing: does Wasserstein actually help before we optimize its cost?
Sinkhorn divergences are differentiable, which matters if optimal transport later graduates
from a diagnostic to a training loss.

\keywords{Sinkhorn divergence, entropic regularization optimal transport, debiased Sinkhorn
divergence, iterative Bregman projections.}

\subsection{OT-5. $W_1$ on increment and gradient PDFs --- \High{}}

\paragraph{What it is.} Rather than transporting the field itself, compute the probability
density of velocity increments $\delta_\ell u$, or of gradients, for both the prediction and
the reference, and compare those one-dimensional distributions using the quantile form of
$W_1$. This is cheap, it has no positivity problem because a PDF is already a probability
measure, and it is sensitive to heavy tails.

\paragraph{Precedent.} Physics-based generative turbulence models are built on exactly these
criteria. A multiscale generative adversarial network uses ``the variance, skewness and
flatness of the increments of the generated field that retrieve respectively the turbulent
energy distribution, energy cascade and intermittency across scales''
\cite{granero2024}. That paper also reports that plain GAN and Wasserstein-GAN baselines
fail to generate fields with turbulent velocity statistics, which gives us a ready-made
negative control.

\rating{High, because this is the cheapest item in the entire tracker that is provably
orthogonal to $L^2$, it requires no mathematical gap to be closed, and it measures
intermittency directly, which is exactly what an over-smoothed emulator destroys first.}

%=========================================================================
\section{Shock geometry and detection}

\subsection{SH-1. Derivative and sharpness detector --- \High{}}

\paragraph{What it is.} A local sensor that flags cells as shocked. The standard forms are
the normalized pressure gradient $|\nabla p| / p$; a compression-gated variant
$\max(0, -\nabla\cdot\mathbf{u}) \cdot |\nabla p| / p$, in which the gate suppresses shear
layers and contact discontinuities; and Mach-gradient sensors. The canonical reference is
Ducros and colleagues, whose stated objective was ``to derive a shock capturing tool able to
treat turbulence with minimum dissipation out of the shock'' for large-eddy simulation
\cite{ducros1999}, which is the same design requirement we have: detect the shock without
contaminating the surrounding turbulence.

\paragraph{Why it matters here.} The meeting observed that a derivative detector ``needs
shock position alignment.'' That is correct: the detector is step one, and step two is
SH-2, which supplies the alignment-free distance between detected sets.

\keywords{shock sensor, Ducros sensor, compression switch, shock-capturing detector,
dilatation--vorticity ratio.}

\paragraph{Pitfall.} Every sensor carries a threshold. Fix the threshold once in the
evaluator and never tune it per model, or the metric becomes gameable.

\subsection{SH-2. Distance between extracted shock sets --- \High{}}

\paragraph{What it is.} Once shock masks or surfaces have been extracted from both fields,
compare them as \emph{geometric objects}. The average symmetric surface distance is the
mean of point-to-surface distances taken in both directions; the 95th-percentile Hausdorff
distance is robust to isolated detector artifacts in a way the raw maximum is not. Taha and
Hanbury survey twenty such measures and select among them systematically
\cite{taha2015}; their treatment also states the specific caveat we need, namely that the
Hausdorff distance ``is generally sensitive to outliers,'' that using it directly is not
recommended when noise and outliers are common, and that the quantile method is the
recommended repair.

\paragraph{This is the direct answer} to ``need to calculate distance between shocks'' from
the Winging It board. Medical imaging solved the problem of deciding whether two thin
surfaces occupy the same place some years ago.

\keywords{average symmetric surface distance, ASSD, 95th percentile Hausdorff distance,
surface distance metrics, boundary F-score.}

\subsection{SH-3. Shock mask overlap --- \Med{}}

Overlap measures such as intersection-over-union and Dice are cheap and easy, but they
degrade badly for \emph{thin} objects: a one-cell shift of a two-cell-thick shock can drive
the intersection-over-union close to zero, which reintroduces the double-penalty problem in
a new costume. The boundary F-score is the thin-structure-aware variant and is the one to
prefer. This is rated Medium rather than High because SH-2 dominates it for our purposes.

\subsection{SH-4. Shock front position and speed error; Rankine--Hugoniot residual --- \Med{}}

Position and speed errors are highly interpretable, but they are well defined only for
isolated, trackable fronts, which is precisely \emph{not} the shocklet case. The
Rankine--Hugoniot residual, which asks whether the jump across a detected front satisfies
the conservation-law jump conditions, is a valuable reference-free physical check, but it
demands careful sampling: never sample states \emph{on} the detected surface, because
captured shocks are smeared across several cells, and instead sample at normal offsets with
tangential averaging.

\rating{Medium: high value in the pure-shock regime, but harder to define for a field
populated by many small shocklets, which is our north-star case.}

%=========================================================================
\section{Norms and function-space metrics}

\subsection{NM-0. Pointwise $L^p$ baselines --- control}
\label{sec:nm0}

\paragraph{What it is.} Mean absolute error, mean squared error, root mean squared error,
and RMSE normalised by the RMS of the reference \emph{fluctuation}. These are not
candidates; they are the control against which the rest of this document is read. A panel
in which $L^2$ scores well everywhere is a broken panel.

\paragraph{Why NRMSE is included.} On the weakly compressible data the density field is
$1.0 \pm 1.8\times10^{-4}$. Normalising by the raw RMS would hide four orders of magnitude
of relative error and would put density and velocity on incomparable scales, so the spatial
mean is removed before normalising.

\paragraph{Response to displacement follows the order of the norm.} For a displacement
$\delta$ small compared with the scale of variation, $f(x+\delta) - f(x) \simeq
\delta\,\partial_x f$, so the two families inherit different powers of $\delta$:
\begin{equation}
\mathrm{MSE} \simeq \delta^{2}\bigl\langle (\partial_x f)^2 \bigr\rangle,
\qquad
\mathrm{MAE} \simeq \delta\,\bigl\langle |\partial_x f| \bigr\rangle .
\label{eq:lp-displacement}
\end{equation}
Measured on $256^2$ vorticity at $\mathrm{Re} = 5\times10^4$ over 21 frames of developed
flow, the ratios per doubling of distance in the sub-cell range are $3.99$, $3.95$, $3.82$
for MSE, against the quadratic prediction of $4$, and $2.00$, $1.98$, $1.93$ for MAE,
against the linear prediction of $2$. Equation~\eqref{eq:lp-displacement} is confirmed to
better than 5\%. Damage below is normalised so that $1$ is the value two statistically
identical but positionally unrelated fields receive.

\begin{center}
\begin{tabular}{@{}lrrrrrrrr@{}}
\toprule
displacement [cells] & 0.125 & 0.25 & 0.5 & 1 & 2 & 4 & 8 & 16 \\
\midrule
MSE   & 0.0004 & 0.0017 & 0.0069 & 0.0262 & 0.0887 & 0.1785 & 0.2789 & 0.4481 \\
MAE   & 0.0221 & 0.0441 & 0.0875 & 0.1692 & 0.3005 & 0.4343 & 0.5835 & 0.6757 \\
NRMSE & 0.0216 & 0.0432 & 0.0858 & 0.1674 & 0.3053 & 0.4392 & 0.5590 & 0.7047 \\
\bottomrule
\end{tabular}
\end{center}

\noindent
MAE therefore assigns 55 times the damage MSE does at an eighth of a cell, and 6.5 times at
one full cell. \emph{If a pointwise metric is wanted that notices sub-cell displacement, MAE
is strictly preferable to MSE}, and the quadratic suppression is why MSE reads as tolerant
of small shifts while being severe at moderate ones. The double penalty does not have a
single onset: its onset depends on the order of the norm. Note also that on this
weakly compressible field $L^2$ saturates over tens of cells rather than immediately,
because the correlation length is long; a shocked field would saturate far faster.

\paragraph{The Gaussian-impostor canary does not catch this family.} Measured, MSE assigns
the spectrum-matched field $0.80$ on vorticity and $0.51$ on velocity: it rejects it firmly,
because a phase-randomised field is pointwise uncorrelated with the reference. $L^p$ metrics
are phase-\emph{sensitive} --- by Plancherel the cross term $\langle fg\rangle$ becomes a sum
of independently randomised phases, so the correlation coefficient goes to zero and the error
approaches its uncorrelated limit. IN-4 is aimed instead at quantities that are functions of
$|\hat f|$ alone, such as BD-1 and BD-2, which score the impostor \emph{perfectly}. This is
recorded here and under IN-4 because the implication is easy to invert.

\paragraph{Redundancy, and a caution.} Across the full degradation ladder the three
baselines correlate at $\rho = 0.995$ (MAE/MSE), $0.970$ (MSE/NRMSE) and $0.968$
(MAE/NRMSE), all above the $0.95$ pruning threshold, so they order the degradations almost
identically. They nonetheless differ by $55\times$ in displacement damage at an eighth of a
cell. A high rank correlation means two metrics \emph{order} damage the same way; it does not
mean they \emph{weight} it the same. For model selection by ranking these are duplicates; as
training losses they are not, and a redundancy matrix read alone would have hidden that.

\rating{No rating. NM-0 is the control, not a candidate. It is listed so that every other
entry has a measured baseline to be read against.}

\subsection{NM-1. Sobolev $H^s$ norm with $s>0$ --- \Low{}}

The meeting has already established the answer here: an $H^1$ penalty helps slightly and
still fails at true discontinuities. The reason is structural, in that a discontinuous
function is not an element of $H^1$ at all, so the gradient term is dominated by a
mesh-resolution artifact rather than by physics. The Low rating reflects not that the idea
is wrong but that we already know the outcome; record it as tested and insufficient rather
than spending further effort.

\subsection{NM-2. Negative Sobolev and dual norms --- \High{}}

\paragraph{What it is.} The interesting direction is the opposite of $H^1$. Negative-order
Sobolev norms are \emph{weaker} than $L^2$: they metrize weak convergence, they remain
finite for rough and discontinuous fields, and, most importantly for us, they are
quantitatively close to the Wasserstein distance. In the periodic case one computes the
homogeneous $\dot H^{-1}$ norm by dividing Fourier coefficients by $|\mathbf{k}|$, so the
cost is one fast Fourier transform.

\paragraph{The relationship to Wasserstein is a theorem, not an analogy.} For infinitesimal
perturbations one has the linearization
\begin{equation}
W_2(\mu, \mu + d\mu) \;=\; \| d\mu \|_{\dot H^{-1}(\mu)} + o(d\mu),
\label{eq:linearization}
\end{equation}
and Peyre integrates this into a non-asymptotic two-sided bound: if $\mu$ and $\nu$ have
densities $f$ and $g$ with $0 < a < f, g < b < \infty$, then
\begin{equation}
b^{-1/2}\, \| f - g \|_{\dot H^{-1}(dx)} \;\le\; W_2(\mu,\nu) \;\le\;
a^{-1/2}\, \| f - g \|_{\dot H^{-1}(dx)}
\label{eq:peyre}
\end{equation}
\cite{peyre2018}. The two distances are therefore equivalent up to constants set by the
density bounds.

\paragraph{The caveat that comes with the theorem.} Equation~\eqref{eq:peyre} degrades when
$a \to 0$, that is when the field has near-vacuum regions. Informally, the equivalence is
good when mass is spread out over the domain and poor when the distribution has regions of
very high and very low density. For our fields this is a real consideration: a density field
with strong rarefactions may sit in the regime where $\dot H^{-1}$ is a weak proxy for
$W_2$, whereas a well-mixed vorticity magnitude field should be comfortably inside it. This
is worth checking empirically on the degradation ladder rather than assuming.

\rating{High: it is nearly free to implement given an FFT, it is well defined for
discontinuous fields where $H^1$ is not, and Equation~\eqref{eq:peyre} makes precise the
sense in which it approximates the metric we actually want. This remains the highest
value-per-line-of-code item in the tracker.}

\keywords{negative Sobolev norm, dual norm misfit, $\dot H^{-1}$ metric, weak convergence
metrization, Benamou--Brenier formula.}

\subsection{NM-3. Mollification split --- \High{}}

\paragraph{What it is.} Convolve the field with a mollifier $\eta_\varepsilon$ at scale
$\varepsilon$ to obtain a smooth part; the remainder $\Delta$ then carries the sharp
structure:
\begin{equation}
f \;=\; \underbrace{f * \eta_\varepsilon}_{\text{smooth part}} \;+\;
\underbrace{\Delta}_{\text{singular remainder}} .
\label{eq:mollify}
\end{equation}
Compare the smooth parts with ordinary $L^2$ or $H^1$ metrics, and compare the remainders
with position-tolerant metrics such as Wasserstein applied to $|\Delta|$, or shock-set
distances applied to the support of $\Delta$.

\paragraph{Why this rates High.} Equation~\eqref{eq:mollify} is not a metric; it is the
\emph{organizing principle} that resolves the whole shocklet problem statement. Shocklets
are simultaneously smooth and sharp, and the split allows the right metric to be applied to
each component instead of demanding that a single metric handle both. Nearly every other
item in this tracker slots naturally into one side or the other.

\paragraph{Prior art to lean on.} This is essentially what Farge's coherent vortex
extraction does in a wavelet basis, catalogued here as BD-3, and what the meteorology
community calls intensity-scale verification, in which forecast skill is decomposed
simultaneously by intensity threshold and by spatial scale using a wavelet decomposition of
the error field \cite{casati2004}.

\paragraph{Open choices, which constitute the mathematical gap.} How should
$\varepsilon$ be chosen: fixed, swept, or tied to the grid? Should the split be iterated
across multiple scales, which would turn it into a wavelet decomposition? Should $\Delta$
be compared as a signed field or through $|\Delta|$?

%=========================================================================
\section{Curvature and differential geometry}

\subsection{CG-1. Level-set and isosurface curvature statistics --- \Med{}}

\paragraph{What it is.} Extract isosurfaces of density, pressure, or vorticity magnitude,
compute their mean and Gaussian curvature, and compare the resulting \emph{distributions}.
Combustion has a mature version of this idea, in which flame-front curvature PDFs are a
standard diagnostic of how turbulence wrinkles a front. For shocked flow the analogue is
shock-front curvature and corrugation statistics.

\keywords{flame front curvature PDF, isosurface curvature turbulence, shock corrugation
spectrum, level-set curvature statistics.}

\rating{Medium: there is real precedent and the implementation cost is moderate, but the
result is sensitive both to the choice of isovalue and to noise in second derivatives.}

\subsection{CG-2. Edge detection --- \Med{}}

Canny and Sobel operators are cheap image-processing shock localizers. Realistically this is
an implementation shortcut for SH-1 rather than an independent idea, so it should be folded
into the shock-detector work item.

\subsection{CG-3. Discrete differential geometry --- \Low{}}

Relevant if we pursue the mesh-comparison path in IN-2, where discrete exterior calculus
supplies structure-preserving operators. As a diagnostic in its own right it is indirect.

\subsection{CG-4. Metric-geometry approach to curvature --- \Low{}}

\paragraph{What it is.} Notions of curvature that require no smooth structure: Alexandrov or
$\mathrm{CAT}(k)$ comparison-triangle curvature, and Ollivier--Ricci curvature, which is
defined through optimal transport between neighborhood measures. Note that Ollivier--Ricci
is built on Wasserstein distance and therefore connects back to OT-1. It has a genuine track
record as a structural descriptor for networks.

\rating{Low for the coming year: we would be inventing the fluid application from scratch,
with no precedent whose failure modes we could borrow, while cheaper items remain untested.
Revisit if we find that we need a mesh-native or graph-native structural descriptor.}

\keywords{Ollivier-Ricci curvature, Forman curvature, Alexandrov curvature, coarse Ricci
curvature networks.}

\subsection{CG-5. Riemann curvature and Arnold's geometric hydrodynamics --- \Low{}}

\paragraph{What it is, now confirmed.} This is the serious reading of ``Riemann curvature''
on the Physics board, and the underlying result is exactly as strong as I claimed. Arnold
extends Euler's theorems on geodesics of $SO(3)$ with a left-invariant metric to arbitrary
Lie groups, in particular to the group $\mathrm{SDiff}(D)$ of volume-preserving
diffeomorphisms; the geodesics of $\mathrm{SDiff}(D)$ are the flows of ideal fluids, which
yields stability criteria for nonlinear hydrodynamics, and he computes the Riemannian
curvature of $\mathrm{SDiff}(D)$ and finds it \emph{negative in most sections}
\cite{arnold1966}. The analytic foundation, showing that these groups can be treated as
infinite-dimensional manifolds on which the Euler equations are genuinely geodesic, is Ebin
and Marsden \cite{ebinmarsden1970}; the modern textbook treatment is Arnold and Khesin
\cite{arnoldkhesin1998}.

\rating{Low as a near-term metric, and I want to be blunt about why. Computing sectional
curvatures of the diffeomorphism group for a discrete velocity field is a research project
in itself, and the physical content it would deliver, namely a predictability horizon
arising from negative curvature, is already available far more cheaply from Lyapunov
exponents and decorrelation times. It is the right answer to a question we can already
answer another way. The verification did not change my rating, but it did confirm that the
idea is real rather than a half-remembered slogan, which is worth knowing.}

%=========================================================================
\section{Vorticity and physics invariants}

\subsection{PH-1. Signed entropy production --- \High{}}

\paragraph{What it is.} In a shocked flow, physical shocks \emph{must} produce entropy;
spurious numerical oscillations produce negative entropy; and over-smeared learned shocks
produce too \emph{little}. Three separate numbers should therefore be reported: an
entropy-violation score capturing unphysical negative production; a signed mismatch against
the reference production; and an excess-dissipation score. A one-sided check detects only
the first of these and will happily pass an over-dissipated mush. The theoretical machinery
for comparing a scheme's entropy production against an entropy-conservative reference flux
is Tadmor's \cite{tadmor1987,tadmor2003}, and entropy production specifically at shocks is
treated by Ismail and Roe \cite{ismailroe2009}.

\rating{High: the diagnostic is reference-free, it is the physically correct admissibility
condition for our regime, and the three-way split directly detects the over-smoothing
failure that we most expect from an ML surrogate.}

\subsection{PH-2. PDE residual in weak form --- \High{}}

This is the ``PDE $\Delta$'' item from the Physics board, done correctly. The
\emph{strong-form} residual, obtained by substituting the field into the differential
operator, is meaningless at a discontinuity, because derivatives of a weak solution are
distributions rather than functions. The \emph{finite-volume weak residual}, consisting of
the cell-average update plus the flux balance minus the source term, remains valid at shocks
and is what should be reported. It is reference-free, cheap, and always available.

\subsection{PH-3. Solver-consistency residual --- \High{}}

\paragraph{What it is.} Take the ML field, advance it by one small step using a trusted
solver, and measure how far the ML model's own next state lies from the solver's. This is
the evaluation-side mirror of solver-in-the-loop training. It requires no ground truth and
remains meaningful long after the trajectory has decorrelated from any reference.

\paragraph{Fits our infrastructure plan.} We are already standing up a Boltzmann solver as
the baseline predictor, per IN-3, so the trusted solver is available at zero marginal cost.

\paragraph{Pitfall.} A small one-step residual can still conceal slow secular drift, so this
should be paired with conservation drift measured over the full rollout.

\subsection{PH-4 and PH-5. Enstrophy and palinstrophy --- \Med{}}

Enstrophy $\int |\boldsymbol{\omega}|^2$ measures small-scale rotational activity, while
palinstrophy $\int |\nabla \boldsymbol{\omega}|^2$ measures the activity of vorticity
gradients and is the natural quantity for the two-dimensional enstrophy cascade. Both are
single scalars and are individually degenerate, in the sense that many wrong fields share a
given enstrophy value, but both are nearly free to compute and make good tripwires. Include
them; do not rely on them.

\paragraph{Enstrophy is implemented, and the measurement sharpens the caveat.} It cannot be
placed on the normalised damage scale at all. That scale is anchored between the reference
value and the value a statistically identical but positionally unrelated field receives ---
and a translated field has \emph{identical} enstrophy by construction, so the two anchors
coincide and the span is zero. The evaluation reports this as ``no dynamic range'' rather
than as a number, which is the honest outcome. Smoothing does make enstrophy fall
monotonically, so it functions exactly as this section describes: a tripwire on a single
field, not a comparison metric between two.

\subsection{PH-6. Helmholtz split --- \High{}}

\paragraph{What it is.} I read the board line ``take curls, but protect from shocks'' as the
observation that vorticity diagnostics are contaminated near shocks, because the velocity
gradient there is dominated by compression rather than rotation. The principled remedy is the
Helmholtz decomposition of the velocity field into a solenoidal, divergence-free, rotational
part and a dilatational, curl-free, compressive part,
\begin{equation}
\mathbf{u} \;=\; \mathbf{u}_{\text{sol}} + \mathbf{u}_{\text{dil}},
\qquad \nabla\cdot\mathbf{u}_{\text{sol}} = 0,
\qquad \nabla\times\mathbf{u}_{\text{dil}} = \mathbf{0},
\end{equation}
after which curl-based diagnostics are computed on $\mathbf{u}_{\text{sol}}$ alone and shock
diagnostics on $\mathbf{u}_{\text{dil}}$. The Ducros-family sensors embody the same logic
pointwise, since they gate on the ratio of dilatation to dilatation plus vorticity in order
to avoid mistaking intense vortical regions for shocks \cite{ducros1999}. Eddy shocklets in
compressible turbulence, our north-star object, were characterized in exactly this
decomposed language \cite{lee1991}.

\rating{High: this is standard practice in the compressible-turbulence community, it is
exactly what the whiteboard line is groping toward, and it yields a principled version of
the mollification split of NM-3 in which the split is defined by physics rather than by an
arbitrary smoothing scale.}

\keywords{Helmholtz decomposition, solenoidal dilatational split, compressible turbulence
decomposition, eddy shocklets.}

%=========================================================================
\section{Pattern and feature detection}

\subsection{PD-1. Hough transform for eddy size --- \Low{} (rating lowered)}
\label{sec:hough}

\paragraph{What it is.} The circle Hough transform detects circular arcs by voting in a
parameter space of centers and radii, while the line version detects straight edges. For our
purposes, circles would correspond to eddies and lines to shock fronts.

\paragraph{What I could not find, stated plainly.} I searched specifically for a
Hough-transform-based eddy-detection precedent and \emph{did not find one}. The vortex- and
eddy-detection literature that does exist is built on streamline geometry, on
velocity-gradient criteria, and on closed-contour methods, not on Hough voting. This does
not prove the idea is bad; it means we would be the first to try it, with no failure modes
to borrow and no baseline to compare against.

\paragraph{One thing worth telling the team anyway.} The Hough transform is essentially a
discretized Radon transform, and the Radon transform is the machinery underlying
\emph{sliced} distances, catalogued as OT-3, and ``project onto planes and compare,''
catalogued as IN-1. Three whiteboard items are the same mathematical object viewed from
different angles.

\rating{Lowered from Medium to Low. The expected precedent is absent, the parameter-space
output is not naturally a metric since one must still define a distance between two Hough
accumulators, and PD-2 covers the same intent with an implemented, documented method. I
would still spend an afternoon visualizing a Hough accumulator of a shocked field, because
the meeting's curiosity about what it looks like is legitimate and cheap to satisfy, but I
would not queue it as a metric candidate ahead of PD-2.}

\subsection{PD-2. Winding angle and Poincar\'e index --- \High{} (rating raised)}

\paragraph{What it is, and it is a real method.} My reading of ``wake number'' as a winding
quantity is now well supported. Sadarjoen and Post present two methods for detecting vortex
structures in two-dimensional vector fields, both based on streamline geometry; the second
``looks at the curvature of a single streamline and computes a winding angle, a metric of
geometric curvature'' \cite{sadarjoen2000}. The documented advantage is directly relevant to
us: because the winding angle does not depend on the velocity magnitude at a single point,
the method \emph{detects weak vortices}, which amplitude-based criteria miss. The documented
disadvantage is equally concrete: a large number of streamlines must be seeded and computed
to maintain complete coverage of the domain. An earlier version of the same geometric
approach appears in \cite{sadarjoen1999}.

\paragraph{The broader family.} The integer-valued relative, the Poincar\'e index of critical
points, underpins vector field topology, introduced for fluid flow data by Helman and
Hesselink \cite{helman1989,helman1991}. Their own framing confirms the dimensional caveat I
raised: the two-dimensional theory of critical points is presented as the \emph{basis} from
which three-dimensional separated flows are then examined, which is to say 2D is where the
method is clean.

\paragraph{The meeting recorded that this was hard to find in the literature.} That was a
keyword problem, as suspected. The work lives in the scientific visualization community
under ``vector field topology'' and ``winding angle,'' not in fluid dynamics journals. A
state-of-the-art survey with a full bibliography is \cite{laramee2007}.

\rating{Raised from Medium to High. An integer- or angle-valued, translation-robust
structural count is genuinely orthogonal to every continuous metric on this list, the method
is implemented and surveyed, its strengths and weaknesses are documented rather than
speculative, and its specific strength (detecting weak vortices independent of local
velocity magnitude) is precisely what an over-smoothed emulator would erase. Start in 2D.}

\subsection{PD-3. Eddy count distribution --- \Med{}}

Once PD-2 or PD-4 has identified eddies, the distributions of their number, size, and
lifetime become natural targets, and those distributions can be compared using the $W_1$
machinery of OT-5. Sadarjoen and Post's title is worth noting here: their method covers
detection, \emph{quantification}, and \emph{tracking} \cite{sadarjoen2000}, so the census
and its time evolution come from the same machinery. The related classical tool is Ripley's
$K$ function, or the pair-correlation function, from spatial point-process statistics, which
is also the rigorous version of the board's ``PDF of pairwise distances'' item, PS-5.

\subsection{PD-4. $Q$-criterion, $\lambda_2$, and Okubo--Weiss --- \Med{}}

These are the established vortex-identification criteria; $\lambda_2$ in particular is due
to Jeong and Hussain \cite{jeonghussain1995}. They are cheap, familiar, and provide a
sensible baseline against which to judge whether PD-2 adds anything. Their known weakness is
that all of them are threshold dependent and none is frame invariant.

%=========================================================================
\section{Topological data analysis}

The meeting described persistence diagrams as low-hanging fruit. I agree, with one important
qualification about which variant.

\subsection{TD-1. One-parameter persistence diagrams --- \High{}}

\paragraph{What it is.} Sweep a threshold through a scalar field such as vorticity magnitude
or density. Structures are born and die as the threshold moves, and each is recorded as a
point $(\text{birth}, \text{death})$ in the persistence diagram. Long-lived features lie far
from the diagonal while noise hugs it. Two diagrams are then compared using the bottleneck
distance or a Wasserstein distance defined on diagrams.

\paragraph{Why it is genuinely low-hanging.} The stability theorem guarantees that a small
perturbation of the field produces a correspondingly small change in the diagram
\cite{cohensteiner2007}, so the metric cannot blow up in the way $L^2$ does under a spatial
shift. Off-the-shelf software exists, including GUDHI, Ripser, giotto-tda, and TTK, so this
is a pure implementation task.

\paragraph{The weather precedent the team remembered is real, and it is from LBNL.}
Muszynski and colleagues recognize atmospheric rivers in climate data using topological data
analysis and machine learning \cite{muszynski2019}. Two details are worth noting. First, the
method is explicitly \emph{threshold-free}, and the authors argue on that basis that it ``may
be useful in evaluating model biases'' in atmospheric-river statistics, which is our use case
almost verbatim. Second, the ground-truth labels were produced by TECA, and the
acknowledgements thank Dmitriy Morozov and Burlen Loring of LBNL's Computational Research
Division for sharing their expertise on computational mathematics. There are people nearby
who have done this.

\rating{High: stable theory, mature software, direct institutional precedent, and
sensitivity to structure in a way that no norm possesses.}

\paragraph{Pitfall.} A persistence diagram is invariant to a great deal, including where in
the domain the features actually sit. Similar topology does not imply correct amplitudes,
shock speeds, or energy transfer, so this belongs in a panel rather than standing alone.

\subsection{TD-2. Two-parameter persistence --- \Low{}}

\paragraph{What it is.} Filtering by two parameters simultaneously, for instance function
value together with spatial scale or distance, which is what the board's note about the
``interplay between distance and function value'' is asking for. The motivation is sound.

\paragraph{Why I nonetheless rate it Low for now.} In two or more parameters there is no
complete discrete invariant analogous to the persistence diagram, which is a theorem-level
obstruction rather than a software gap. Tooling has improved since the situation Lesnick and
Wright described in 2015, when they noted that there was ``no publicly available software
which extends the usual persistent homology pipeline for exploratory data analysis to the
multidimensional persistence''; their RIVET, the Rank Invariant Visualization and Exploration
Tool, was built to fill that gap and provides interactive visualization of the barcodes of
one-dimensional affine slices of a two-dimensional persistence module
\cite{lesnickwright2015}, with efficient presentation and Betti-number algorithms added
later \cite{lesnickwright2022}. It is usable but remains research-grade relative to the
one-parameter tooling.

\paragraph{Recommendation.} Park it. If TD-1 succeeds and its principal failure turns out to
be scale confounding, revisit this item then, starting from RIVET.

\subsection{TD-3. Well groups --- \Low{}}

Well groups measure the robustness of level sets to perturbation. The idea is elegant and
the computational machinery is less mature than that of persistence. The recommendation is
the same as for TD-2.

\subsection{TD-4. Reeb graphs and Morse--Smale complexes --- \Med{}}

These are skeletons of the level-set connectivity and of the gradient flow respectively, and
both are implemented in TTK \cite{tierny2018}. They are more informative than a persistence
diagram about \emph{where} structures are located, but the metrics used to compare such
graphs, including interleaving and edit distances, are less standardized.

\subsection{TD-5. Jacobi sets --- \Med{}}

\paragraph{Worth calling out specifically, and now confirmed to be exactly what I hoped.}
The Jacobi set of two Morse functions defined on a common $d$-manifold is the set of critical
points of the restrictions of one function to the level sets of the other; equivalently, and
more usefully for us, it is ``the set of points where the gradients of the functions are
parallel'' \cite{edelsbrunnerharer2004}. That makes it natively a \emph{comparison} object,
which is unusual, since most topological summaries characterize a single field and then
compare the summaries. For prediction against ground truth this is structurally the right
shape of tool. Two further points from the source strengthen the case: there is a
polynomial-time algorithm for the piecewise-linear analogue on triangulations, so it is
computable; and the authors' own motivating example is oceanography, where researchers study
several attributes of water on a common domain, which is close to our setting.

\rating{Medium, but with the highest upside per unit of obscurity in the whole topology
theme. The gap classification improves from ``math and implementation'' to
``implementation'' now that the polynomial-time PL algorithm is confirmed. If someone wants a
distinctive contribution, this is where I would look.}

%=========================================================================
\section{Basis decompositions and the turbulence transform}

\subsection{BD-1 and BD-2. Fourier spectrum and two-point correlation --- \Med{}}

Both are standard and cheap, and we should compute them, provided we understand what they
are \emph{not}. The two-point correlation and the power spectrum are Fourier transforms of
one another, so they carry the same information and running both does not constitute two
independent checks. More importantly, both discard all Fourier \emph{phase}, which is where
structure resides, so matching a spectrum is a weak constraint. Include them as baselines
and do not treat them as evidence of structural fidelity.

\subsection{BD-3. Wavelets and coherent vortex extraction --- \High{}}

\paragraph{What it is.} Wavelets are localized in both space and scale, so they do not smear
a shock across all wavenumbers the way Fourier modes do. As Farge puts it in the review that
introduced this toolset to the turbulence community, wavelets ``unfold a signal, or a field,
into both space and scale,'' and their limited spatial support means the analysis ``can be
performed locally on the signal, as opposed to the Fourier transform which is inherently
nonlocal due to the space-filling nature of the trigonometric functions'' \cite{farge1992}.
Coherent vortex extraction then splits a turbulent field into a coherent part, carried by a
small number of large wavelet coefficients, and an incoherent, Gaussian-like background
carried by the many small ones \cite{fargeschneider1999,farge2001}. A more recent survey of
wavelet methods in CFD is \cite{schneidervasilyev2010}.

\paragraph{Note the coincidence with your own whiteboard.} Coherent vortex extraction is
Equation~\eqref{eq:mollify} carried out in a principled, shift-localized basis, with a
data-driven rather than arbitrary split. If the mollification idea is the right organizing
principle, then wavelets are its mature implementation. The Farge quotation above is also, in
effect, the argument for why wavelets should beat Fourier for shocklets specifically.

\keywords{coherent vortex extraction, CVS, wavelet thresholding turbulence, local
intermittency measure, wavelet methods computational fluid dynamics.}

\subsection{BD-4. Wavelet scattering and phase harmonics --- \High{}}

\paragraph{What it is.} This is my candidate for the board's ``some kind of turbulence
transform,'' and the verification supports it. A scattering transform cascades wavelet
transforms with modulus nonlinearities. The first layer depends on a single scale and
orientation; the second layer probes the \emph{coupling between two oriented scales}, which
is exactly the information a power spectrum discards. It is described in the astrophysics
literature as ``a low-variance statistical description of non-Gaussian processes'' that is
``inspired by convolutional neural networks but that does not require any training stage''
\cite{regaldo2020}, originating with Mallat \cite{mallat2012} and developed by Bruna and
Mallat \cite{brunamallat2013}.

\paragraph{Why this fits our problem specifically.} First, the transform is translation
invariant and stable to deformations, which is the formal version of the requirement that a
slightly shifted shock should not blow up the metric. Second, its second-layer coefficients
measure cross-scale coupling, which is what distinguishes real turbulence from a Gaussian
field with an identical spectrum. Third, it has been used on magnetohydrodynamic turbulence
simulations and observations in a field that, like ours, never has phase-aligned ground
truth: Allys and colleagues introduced the reduced wavelet scattering transform for the
interstellar medium \cite{allys2019}, and the approach was subsequently extended to
polarized dust emission \cite{regaldo2020} and used to classify MHD simulations
\cite{saydjari2021}.

\keywords{wavelet scattering transform, reduced wavelet scattering transform, RWST,
scattering covariance, wavelet phase harmonics, Mallat scattering.}

\rating{High: this is the only item on the whiteboard that plausibly delivers a single
descriptor family sensitive to smooth structure, sharp structure, and intermittency
simultaneously, and therefore the only real candidate for a unified shocklet metric. The
astrophysical track record on compressible MHD turbulence is close enough to our regime to be
genuinely informative.}

%=========================================================================
\section{Probabilistic and distributional metrics}

\subsection{PS-1. Energy distance --- \Med{}}

The energy distance is a kernel-family distance between distributions, and is in fact
maximum mean discrepancy with a particular kernel. The board's own caveat, that it ``does
some averaging'' and is therefore hard for heavy tails, is the correct objection and is
exactly why Wasserstein is preferred here. Keep it as a cheap baseline that demonstrates the
difference, and do not expect it to win.

\subsection{PS-2. CRPS --- \High{}, conditional on ensembles}

The continuous ranked probability score is the strictly proper scoring rule on which the
weather community trains \cite{gneitingraftery2007}. It reduces to absolute error for a
deterministic forecast, so \emph{it conveys nothing beyond mean absolute error unless an
ensemble or a generative model is available}. This should be flagged clearly in the plan: if
the Boltzmann-degradation study produces single deterministic fields, then PS-2 and PS-3 are
not applicable and should be marked N/A rather than left as untested.

\subsection{PS-3. Spread-to-skill ratio --- \High{}, conditional on ensembles}

\paragraph{What it is, with a computational gotcha worth knowing in advance.} This is the
ratio of ensemble spread to the root-mean-square error of the ensemble mean, and for a
well-calibrated ensemble it should sit near unity. Fortin and colleagues justify the practice
on the basis of exchangeability between forecasts and observations, and, more practically,
they show that the average spread must be computed as \emph{the square root of the average
ensemble variance} rather than as the average of the spreads \cite{fortin2014}. Using the
naive averaging gives a downward-biased spread and therefore a spurious diagnosis of
underdispersion; in their operational example, correcting the estimator changed the verdict
from ``underdispersed'' to ``excellent agreement, with possible overdispersion at long lead
times.'' If we implement this from scratch we should get the estimator right the first time.

\subsection{PS-4. Heavy-tail quantifiers --- \High{}}

There are three tiers here, in increasing order of sophistication. First, the flatness or
kurtosis of velocity increments and gradients, which for a Gaussian field equals exactly
three and which for real turbulence greatly exceeds three and grows toward small scales,
giving a one-number Gaussianity alarm. Second, direct comparison of tail exceedance
probabilities $P(|q| > a)$. Third, extreme-value-theory tail-index estimation via the Hill
estimator or a peaks-over-threshold fit to a generalized Pareto distribution. The first tier
is precisely one of the three physics-based criteria used to train and evaluate the
turbulence GAN of \cite{granero2024}.

\rating{High: intermittency is the first property an over-smoothed emulator destroys and the
last one an $L^2$ loss notices, and the first tier costs almost nothing.}

\keywords{flatness kurtosis velocity increments, intermittency exponent, Hill estimator tail
index, peaks over threshold generalized Pareto.}

\subsection{PS-5. PDF of pairwise distances --- \Med{}}

The board's ``PDF of pairwise distances'' formalizes as the pair-correlation function, or
Ripley's $K$ function, from spatial point-process statistics, once point-like features such
as eddy centers or shocklet locations have been identified. This item therefore depends on
PD-2 or PD-3 landing first.

%=========================================================================
\section{Imports from weather verification}

The meeting's instinct that the weather community has already solved parts of this problem
is correct. The verification turned up something better than a vocabulary list: an explicit
statement of our problem.

\begin{description}[leftmargin=*,style=nextline]
\item[The double penalty problem, stated in the literature.] Wernli and colleagues write
that ``the classical example to illustrate the limitations of gridpoint-based error measures
is the `double penalty problem': a prediction of a precipitation structure that is correct in
terms of amplitude, size, and timing but (maybe only slightly) incorrect concerning position
is very poorly rated'' \cite{wernli2008}. Replace ``precipitation structure'' with ``shock''
and that is the first bullet of our meeting notes, published in 2008. Searching this term
surfaces two decades of work on the pathology.
\item[SAL.] The structure--amplitude--location measure decomposes forecast quality into three
components, where the location component combines the displacement of the field's center of
mass with the error in the weighted-average distance of features from that center
\cite{wernli2008}. The decomposition philosophy is the same as our NM-3 split.
\item[MODE.] Object-based verification, which uses convolution, smoothing, and thresholding
to define meaningful objects and then a matching algorithm to pair them between forecast and
observation \cite{davis2006a,davis2006b,davis2009}. This is the closest existing analogue to
``extract shocklets from both fields, then compare the sets.''
\item[CRA.] Contiguous rain area verification, which decomposes error into displacement,
volume, and pattern components \cite{ebertmcbride2000}.
\item[Intensity-scale verification.] A wavelet-based approach that scores skill jointly by
intensity threshold and spatial scale \cite{casati2004}; the direct ancestor of BD-3 applied
to verification rather than to analysis.
\item[Threshold-free topological detection.] Muszynski et al.\ \cite{muszynski2019},
discussed under TD-1.
\item[The strategic lesson.] The meeting already drew this, and I would put it in the
repository README: the diagnostic catalog and the training loss are separate artifacts with
separate requirements. A diagnostic may be non-differentiable, expensive, and discontinuous;
a loss may not.
\end{description}

%=========================================================================
\section{Infrastructure and protocol}

\subsection{IN-1. Project onto planes and compare --- \Med{}}

As noted under OT-3 and PD-1, this is the Radon and slicing idea, and it is largely subsumed
by sliced Wasserstein. Retain it as a \emph{visualization} task rather than as a separate
metric.

\subsection{IN-2. Adaptive meshes and cross-mesh comparison --- \Med{}}

Two distinct questions hide inside this item. The first is plumbing: how do we compare fields
that live on different meshes at all? The answer is to remap both conservatively onto a
common analysis grid, to compare cell averages rather than point samples, and to record the
remapping operator as part of the metric definition. The second is diagnostic: an adaptive
solver's refinement indicator is a statement about where the solver believes structure lives,
so comparing refinement maps compares structural beliefs. The second question is genuinely
interesting, while the first is necessary infrastructure for everything else and should be
scheduled early even though it is not itself a metric.

\subsection{IN-3. Degradation ladder and acceptance test --- \High{}}

This turns the meeting's plan into a falsifiable procedure. The protocol below is my
synthesis rather than something drawn from the literature.

\begin{enumerate}[leftmargin=*]
\item \textbf{Reference.} Direct numerical simulation at $\mathrm{Re} \approx 500$, isotropic
and doubly periodic shear, as agreed in the meeting.
\item \textbf{Predictor ladder.} The Boltzmann solver degraded in controlled steps through
grid coarsening, and ideally along a second, \emph{orthogonal} degradation axis such as added
noise or artificial viscosity, so that we can distinguish metrics that respond to resolution
specifically from metrics that respond to badness generally.
\item \textbf{Computation.} Evaluate every candidate metric at every rung of the ladder.
\item \textbf{Acceptance criterion.} A metric is admitted to the panel if it is
\emph{monotone in degradation with a high Spearman rank correlation}. Non-monotone metrics
are actively dangerous, because model selection performed against them can move us in the
wrong direction.
\item \textbf{Sensitivity.} Record where each metric first departs from its clean value
relative to the point at which a human observer sees visible degradation. A metric that fires
only after the field already looks obviously wrong is not earning its place in the panel.
\item \textbf{Cost.} Measure wall-clock time per evaluation, so that we can construct the
cost-accuracy frontier the meeting asked for.
\item \textbf{Redundancy.} Correlate metrics against one another across the ladder and prune
near-duplicates. The goal is an \emph{orthogonal} panel, not a long one.
\end{enumerate}

\subsection{IN-4. Gaussian-impostor canary --- \High{}}

\paragraph{What it is.} Add one further ``predictor'' to the ladder that is not a degraded
solver at all, namely a Gaussian random field constructed to match the reference energy
spectrum exactly. Such a field has the correct spectrum and the correct two-point
correlation, but it has zero skewness, a vanishing third-order structure function, zero
energy flux, a flatness of exactly three, and no phase coupling whatsoever. It is not
turbulence in any physical sense.

\paragraph{Implemented, with one correction worth recording.} The phases must be taken from
the transform of a \emph{real} Gaussian field rather than drawn as $\exp(i\,U(0,2\pi))$
directly. Hand-drawn phases violate Hermitian symmetry at the self-conjugate modes --- $k=0$,
the Nyquist row and column, and the Nyquist corners --- so the inverse transform silently
discards the imaginary part and corrupts the very amplitude spectrum the construction is
meant to preserve. A first implementation that way produced a flatness of $47.9$; the
corrected version gives a bit-identical amplitude spectrum, a two-point correlation matching
to relative $10^{-17}$, and a flatness of $3.01$ against a reference flatness of $17.06$.

\paragraph{Measured: it does not catch the $L^p$ family.} MSE assigns the impostor $0.80$ on
vorticity and $0.51$ on velocity, where $1$ is the value two unrelated fields receive --- it
rejects it firmly. That is the correct outcome and not a failure of the canary: $L^p$ metrics
are phase-\emph{sensitive}, so randomising the phases is the most damaging thing that can be
done to them. The canary's discriminating power is one-sided. It can only reject, and a
metric passing it conveys almost nothing. The quantities it does catch are those that are
functions of $|\hat f|$ alone --- the energy spectrum (BD-1), the two-point correlation
(BD-2), second-order structure functions --- each of which scores the impostor perfectly. A
report in which every implemented metric passes IN-4 should therefore not be read as
reassuring until such a metric is in the panel. See also Section~\ref{sec:nm0}.

\paragraph{It is one coordinate of two.} A complementary probe is the reference translated by
a large random offset: on a periodic domain this preserves every single- and multi-point
statistic \emph{exactly} (measured variance ratio $1.000000$, flatness matching to three
decimals) while destroying pointwise alignment. The impostor asks whether a metric sees
beyond second-order statistics; the translated field asks whether it tolerates displacement.
Together they place a metric on a two-way map whose useful quadrant --- displacement-tolerant
and still structure-sensitive --- is what OT-1, TD-1 and BD-4 each claim to occupy. The
translated field also serves as the measured anchor that normalises every damage score.

\paragraph{Why it belongs in the protocol.} Any metric that ranks the Gaussian impostor as
\emph{good} is thereby confirmed to be phase-blind and must not be used alone for model
selection. This is a cheap and decisive discriminator that separates the genuinely structural
metrics, such as TD-1, BD-4, PS-4, and PH-6, from those that see only second-order
statistics, such as BD-1 and BD-2. It costs roughly one afternoon and should immediately
halve the candidate list.

%=========================================================================
\section{Recommended first wave}

If we can complete only six items before the next meeting, I would choose the following,
ordered roughly by value per unit of effort.

\begin{enumerate}[leftmargin=*]
\item \textbf{IN-4, the Gaussian-impostor canary.} The cheapest possible experiment, it
prunes the candidate list immediately and requires no new machinery.
\item \textbf{NM-2, the negative Sobolev norm $\dot H^{-1}$.} An FFT and a division by
$|\mathbf{k}|$, with Equation~\eqref{eq:peyre} telling us exactly what we are approximating
and under what condition the approximation is good.
\item \textbf{TD-1, sublevel-set persistence on vorticity and density.} Off-the-shelf
software, stable theory, local institutional precedent, and it is already the team's stated
next step.
\item \textbf{OT-5 together with PS-4, increment-PDF Wasserstein and flatness.} No positivity
problem to solve, direct measurement of intermittency, and roughly one afternoon of work.
\item \textbf{SH-1 together with SH-2, the shock detector plus surface distances.} This
closes the ``need to calculate distance between shocks'' gap using imported medical-imaging
metrics.
\item \textbf{PH-2 together with PH-3, the weak PDE residual and the solver-consistency
residual.} Both are reference-free, and PH-3 is nearly free once the Boltzmann solver is in
the repository.
\end{enumerate}

\noindent
For the second wave I would propose BD-4, wavelet scattering and phase harmonics; PH-6, the
Helmholtz split; PD-2, the winding angle in two dimensions, which has moved up on the
strength of the verified precedent; and OT-1 with OT-2, full Wasserstein together with the
signed-field fix.

%=========================================================================
\section{Implementation status}
\label{sec:implementation}

The evaluation harness exists and runs on the real data. This section records what is built
and, more usefully, the corrections that building it forced --- each was found by measurement
and each would otherwise have corrupted results quietly.

\subsection{What exists}

\begin{description}[leftmargin=*,style=nextline]
\item[Three plugin registries.] Metrics, degradations, and report renderers are each
extended by writing one decorated function; discovery is by name from configuration, so
nothing has an import list to maintain. Data readers are a deliberately closed set.
\item[The ladder (IN-3).] Twenty degradation operators in seven families: kernel smoothing
(Gaussian, box, median, disk, Epanechnikov), spectral filters (ideal and Butterworth
low- and high-pass, band attenuation), translations (whole-cell and sub-cell), coarsening
and its non-conservative control, noise, and pointwise gain and bias. Acceptance quantities
are computed per axis: rank correlation with a moving-block bootstrap interval, the fraction
of frames strictly ordered, adjacent-rung separability, and the rungs at which a metric first
departs from clean and saturates.
\item[The analysis grid (IN-2).] Conservative block averaging onto a common grid, with the
remapping operator recorded alongside every number.
\item[The canary (IN-4).] Implemented and measured; see the corrected construction above.
\item[Reporting.] Each run produces a self-contained folder holding the raw numbers, the
resolved configuration, full provenance, and a \LaTeX{} document that compiles standalone and
renders on Overleaf.
\end{description}

\subsection{Five corrections the measurements forced}

\begin{enumerate}[leftmargin=*]
\item \textbf{Rank correlation must be computed per frame, not pooled over frames.} The
density perturbation grows six orders of magnitude along the trajectory, so the worst rung
early is numerically smaller than the mildest rung late. Pooling therefore measures the
flow's own evolution rather than the metric's response: every density axis was perfectly
ordered \emph{within} every frame while the pooled value read between $0.10$ and $0.91$
depending on the axis.

\item \textbf{Derived fields must be recomputed after a remap, never averaged.}
Block-averaging vorticity yields a field that is not the curl of the velocity beside it. The
difference is $5.6\%$, $18.3\%$ and $25.9\%$ at coarsening factors $2$, $4$ and $8$.
Relatedly, a stored vorticity must not be mixed with a recomputed one: the solver's lattice
stencil and a spectral derivative differ by $8.1\%$ rms (correlation $0.997$), so using one
at the native grid and the other at a coarse grid would make a grid-independence comparison
measure the discretisation rather than the grid.

\item \textbf{High-pass filters must preserve the spatial mean.} Deleting the $k=0$ mode of
a density field removes a component four orders of magnitude larger than anything the cutoff
controls. Before this was corrected every rung returned an identical damage of $2.7\times
10^{7}$ and the axis carried no ordering whatsoever.

\item \textbf{The unrelated-field anchor must be measured, not inferred.} Substituting the
largest configured translation understates it --- a 16-cell displacement reaches only about
$0.6$ of the true value here --- which inflates every damage score by roughly $1.6\times$. A
distant frame of the same trajectory is also wrong, and for a reason worth recording: this
flow decays, so a frame $4000$ steps away has $0.70$ of the variance and a different
flatness, and it scores \emph{closer} to the reference than a true statistical twin does
purely because the field has weakened.

\item \textbf{Severity ranges must follow each field's spectrum.} Fixed cutoffs applied to
every field alike produce flags that point at the axis rather than at the metric. On
vorticity, high-pass damage runs $0.566$, $0.503$, $0.548$, $0.580$ across cutoffs $2$, $4$,
$8$, $16$ --- non-monotone and indistinguishable, because a cutoff of $2$ already removes
almost all the energy. This is the one correction not yet applied.
\end{enumerate}

\subsection{Consequences for the protocol as written}

Two items in Section~\ref{sec:master} cannot be executed exactly as specified with the data
currently available, and both are recorded rather than worked around. The protocol calls for
a reference at $\mathrm{Re} \approx 500$; every doubly periodic dataset on disk is
$\mathrm{Re} = 5\times10^4$ to $10^6$. And the north-star shocklet case is not what the
present target provides --- \texttt{weakly\_compressible\_isoT} at $\mathrm{Ma} = 0.1$ has
density varying by under $4\%$ and no shocks --- so SH-1, SH-2 and PH-1 have nothing valid to
be validated against yet. Shocked datasets are scheduled; the harness is format-agnostic and
will run on them unchanged.

%=========================================================================
\section{Confidence}

\paragraph{Verified against primary sources.}
Every citation in the bibliography of Section~\ref{sec:bib} has been checked. The items that
most directly support ratings are: the Wasserstein-for-shifted-signals argument and its
$s^2$ convexity statement \cite{engquist2014,engquistyang,yang2016}; graph-space optimal
transport as the fix for signed, non-normalized data
\cite{metivier2016gji,metivier2018geophysics,metivier2018eage,metivier2019ip}; the two-sided
$W_2$ versus $\dot H^{-1}$ bound with its density conditions \cite{peyre2018}; the
double-penalty statement and the SAL decomposition \cite{wernli2008}; object-based
verification \cite{davis2006a,davis2006b,davis2009}, CRA \cite{ebertmcbride2000}, and
intensity-scale verification \cite{casati2004}; segmentation surface metrics with the
Hausdorff outlier caveat \cite{taha2015}; the Ducros shock sensor and its design goal
\cite{ducros1999}; the winding-angle vortex detector with its stated advantage for weak
vortices and its streamline-seeding cost \cite{sadarjoen2000,sadarjoen1999}, plus vector
field topology \cite{helman1989,helman1991,laramee2007}; persistence stability
\cite{cohensteiner2007}, the threshold-free atmospheric-river application with its
model-bias framing and LBNL acknowledgements \cite{muszynski2019}, RIVET and the
two-parameter tooling situation \cite{lesnickwright2015,lesnickwright2022}, Jacobi sets as
the set where gradients are parallel together with the polynomial-time PL algorithm
\cite{edelsbrunnerharer2004}, and TTK \cite{tierny2018}; wavelets and coherent vortex
extraction \cite{farge1992,fargeschneider1999,farge2001,schneidervasilyev2010}; the wavelet
scattering transform and its astrophysical applications
\cite{mallat2012,brunamallat2013,allys2019,regaldo2020,saydjari2021}; the spread-versus-RMSE
estimator subtlety \cite{fortin2014}; proper scoring rules \cite{gneitingraftery2007}; the
physics-based turbulence GAN criteria and its failed baselines \cite{granero2024}; Arnold's
negative sectional curvature result and its analytic foundation
\cite{arnold1966,ebinmarsden1970,arnoldkhesin1998}; entropy-stability theory
\cite{tadmor1987,tadmor2003,ismailroe2009}; eddy shocklets \cite{lee1991}; and $\lambda_2$
\cite{jeonghussain1995}.

\paragraph{Not substantiated.}
I searched specifically for a Hough-transform-based eddy-detection precedent and found none.
PD-1 is therefore recorded as an untested original idea rather than as an import, and its
rating was lowered accordingly. This is the only claim from the previous revision that did
not survive checking.

\paragraph{My judgment rather than literature consensus.}
Every promise rating and its accompanying justification, including the two changes made in
this revision. The claim that the Hough transform, the Radon transform, and slicing collapse
three whiteboard items into one work item. The claim that the Helmholtz split is the
principled reading of ``take curls but protect from shocks.'' The reading of ``wake number''
as a winding quantity, which the discovery of the winding-angle method supports but does not
prove. The identification of wavelet scattering as the best candidate for ``some kind of
turbulence transform.'' The entire IN-3 protocol, including the Spearman acceptance criterion
and the orthogonal second degradation axis. The IN-4 canary proposal. The first-wave
ordering. The observation that Equation~\eqref{eq:peyre} degrades for fields with
near-vacuum regions is a straightforward reading of the stated hypotheses, but the
implication for our specific density fields is my inference and should be checked
empirically.

%=========================================================================
\section{Connections worth keeping in view}

Several themes on the boards are the same mathematics wearing different clothes, and
recognizing this reduces the total work.

\begin{itemize}[leftmargin=*]
\item \textbf{Hough $\leftrightarrow$ Radon $\leftrightarrow$ sliced Wasserstein
$\leftrightarrow$ ``project onto planes.''} One transform, four whiteboard entries.
\item \textbf{Mollification split $\leftrightarrow$ coherent vortex extraction
$\leftrightarrow$ Helmholtz decomposition $\leftrightarrow$ intensity-scale verification.}
One organizing principle, namely separate the smooth from the sharp and metricize each
appropriately, with four communities' implementations.
\item \textbf{Wasserstein $\leftrightarrow$ $\dot H^{-1}$ $\leftrightarrow$ Ollivier--Ricci
curvature.} All are built on transport; the first two are quantitatively equivalent under the
density conditions of Equation~\eqref{eq:peyre}, which makes $\dot H^{-1}$ the cheap proxy.
\item \textbf{Double penalty (weather) $\leftrightarrow$ cycle skipping (seismic)
$\leftrightarrow$ our shock-shift problem.} One pathology, three names, two of which carry a
mature literature we can import wholesale.
\item \textbf{Persistence diagrams $\leftrightarrow$ eddy census $\leftrightarrow$ Ripley's
$K$ $\leftrightarrow$ MODE object matching.} All are instances of ``identify features, then
compare their distributions,'' differing only in how the features are defined and whether the
features are matched pairwise before comparison.
\end{itemize}

%=========================================================================
\section{Bibliography}
\label{sec:bib}

\begin{thebibliography}{99}
\small

\bibitem{allys2019} Allys, E., Levrier, F., Zhang, S., Colling, C., Regaldo-Saint Blancard,
B., Boulanger, F., Hennebelle, P. \& Mallat, S. (2019). The RWST, a comprehensive statistical
description of the non-Gaussian structures in the ISM. \emph{Astronomy \& Astrophysics} 629,
A115. \url{https://arxiv.org/abs/1905.01372}

\bibitem{arnold1966} Arnold, V.\ I. (1966). Sur la g\'eom\'etrie diff\'erentielle des groupes
de Lie de dimension infinie et ses applications \`a l'hydrodynamique des fluides parfaits.
\emph{Annales de l'Institut Fourier} 16(1), 319--361.
\url{https://doi.org/10.5802/aif.233}

\bibitem{arnoldkhesin1998} Arnold, V.\ I. \& Khesin, B.\ A. (1998). \emph{Topological Methods
in Hydrodynamics}. Applied Mathematical Sciences 125, Springer, New York.

\bibitem{brunamallat2013} Bruna, J. \& Mallat, S. (2013). Invariant scattering convolution
networks. \emph{IEEE Transactions on Pattern Analysis and Machine Intelligence} 35(8),
1872--1886. \url{https://doi.org/10.1109/TPAMI.2012.230}

\bibitem{casati2004} Casati, B., Ross, G. \& Stephenson, D.\ B. (2004). A new intensity-scale
approach for the verification of spatial precipitation forecasts.
\emph{Meteorological Applications} 11(2), 141--154.
\url{https://doi.org/10.1017/S1350482704001239}

\bibitem{cohensteiner2007} Cohen-Steiner, D., Edelsbrunner, H. \& Harer, J. (2007). Stability
of persistence diagrams. \emph{Discrete \& Computational Geometry} 37(1), 103--120.
\url{https://doi.org/10.1007/s00454-006-1276-5}

\bibitem{davis2006a} Davis, C., Brown, B. \& Bullock, R. (2006). Object-based verification of
precipitation forecasts. Part I: Methodology and application to mesoscale rain areas.
\emph{Monthly Weather Review} 134(7), 1772--1784.
\url{https://doi.org/10.1175/MWR3145.1}

\bibitem{davis2006b} Davis, C., Brown, B. \& Bullock, R. (2006). Object-based verification of
precipitation forecasts. Part II: Application to convective rain systems.
\emph{Monthly Weather Review} 134(7), 1785--1795.
\url{https://doi.org/10.1175/MWR3146.1}

\bibitem{davis2009} Davis, C.\ A., Brown, B.\ G., Bullock, R. \& Halley Gotway, J. (2009).
The Method for Object-Based Diagnostic Evaluation (MODE) applied to numerical forecasts from
the 2005 NSSL/SPC Spring Program. \emph{Weather and Forecasting} 24(5), 1252--1267.
\url{https://doi.org/10.1175/2009WAF2222241.1}

\bibitem{ducros1999} Ducros, F., Ferrand, V., Nicoud, F., Weber, C., Darracq, D., Gacherieu,
C. \& Poinsot, T. (1999). Large-eddy simulation of the shock/turbulence interaction.
\emph{Journal of Computational Physics} 152(2), 517--549.
\url{https://doi.org/10.1006/jcph.1999.6238}

\bibitem{ebertmcbride2000} Ebert, E.\ E. \& McBride, J.\ L. (2000). Verification of
precipitation in weather systems: Determination of systematic errors.
\emph{Journal of Hydrology} 239(1--4), 179--202.
\url{https://doi.org/10.1016/S0022-1694(00)00343-7}

\bibitem{ebinmarsden1970} Ebin, D.\ G. \& Marsden, J.\ E. (1970). Groups of diffeomorphisms
and the motion of an incompressible fluid. \emph{Annals of Mathematics} 92(1), 102--163.
\url{https://doi.org/10.2307/1970699}

\bibitem{edelsbrunnerharer2004} Edelsbrunner, H. \& Harer, J. (2004). Jacobi sets of multiple
Morse functions. In \emph{Foundations of Computational Mathematics, Minneapolis 2002}
(Cucker, F., DeVore, R., Olver, P. \& S\"uli, E., eds.), London Mathematical Society Lecture
Note Series 312, Cambridge University Press, pp.\ 37--57.
\url{https://doi.org/10.1017/CBO9781139106962.003}

\bibitem{engquist2014} Engquist, B. \& Froese, B.\ D. (2014). Application of the Wasserstein
metric to seismic signals. \emph{Communications in Mathematical Sciences} 12(5), 979--988.
\url{https://doi.org/10.4310/CMS.2014.v12.n5.a7}

\bibitem{engquistyang} Engquist, B. \& Yang, Y. Optimal transport based seismic inversion:
Beyond cycle skipping. \emph{Communications on Pure and Applied Mathematics}.
\url{https://yunany.github.io/Publications/11_OT_CPAM.pdf}

\bibitem{farge1992} Farge, M. (1992). Wavelet transforms and their applications to turbulence.
\emph{Annual Review of Fluid Mechanics} 24, 395--458.
\url{https://doi.org/10.1146/annurev.fl.24.010192.002143}

\bibitem{fargeschneider1999} Farge, M., Schneider, K. \& Kevlahan, N. (1999). Non-Gaussianity
and coherent vortex simulation for two-dimensional turbulence using an adaptive orthogonal
wavelet basis. \emph{Physics of Fluids} 11(8), 2187--2201.
\url{https://doi.org/10.1063/1.870080}

\bibitem{farge2001} Farge, M., Pellegrino, G. \& Schneider, K. (2001). Coherent vortex
extraction in 3D turbulent flows using orthogonal wavelets.
\emph{Physical Review Letters} 87(5), 054501.
\url{https://doi.org/10.1103/PhysRevLett.87.054501}

\bibitem{fortin2014} Fortin, V., Abaza, M., Anctil, F. \& Turcotte, R. (2014). Why should
ensemble spread match the RMSE of the ensemble mean? \emph{Journal of Hydrometeorology}
15(4), 1708--1713. \url{https://doi.org/10.1175/JHM-D-14-0008.1}

\bibitem{gneitingraftery2007} Gneiting, T. \& Raftery, A.\ E. (2007). Strictly proper scoring
rules, prediction, and estimation. \emph{Journal of the American Statistical Association}
102(477), 359--378. \url{https://doi.org/10.1198/016214506000001437}

\bibitem{granero2024} Granero-Belinch\'on, C. \& Cabeza Gallucci, M. (2024). A multiscale and
multicriteria generative adversarial network to synthesize 1-dimensional turbulent fields.
\emph{Machine Learning: Science and Technology} 5(2), 025032.
\url{https://doi.org/10.1088/2632-2153/ad43b3}

\bibitem{helman1989} Helman, J.\ L. \& Hesselink, L. (1989). Representation and display of
vector field topology in fluid flow data sets. \emph{Computer} 22(8), 27--36.
\url{https://doi.org/10.1109/2.35197}

\bibitem{helman1991} Helman, J.\ L. \& Hesselink, L. (1991). Visualizing vector field topology
in fluid flows. \emph{IEEE Computer Graphics and Applications} 11(3), 36--46.
\url{https://doi.org/10.1109/38.79452}

\bibitem{ismailroe2009} Ismail, F. \& Roe, P.\ L. (2009). Affordable, entropy-consistent Euler
flux functions II: Entropy production at shocks.
\emph{Journal of Computational Physics} 228(15), 5410--5436.
\url{https://doi.org/10.1016/j.jcp.2009.04.021}

\bibitem{jeonghussain1995} Jeong, J. \& Hussain, F. (1995). On the identification of a vortex.
\emph{Journal of Fluid Mechanics} 285, 69--94.
\url{https://doi.org/10.1017/S0022112095000462}

\bibitem{laramee2007} Laramee, R.\ S., Hauser, H., Zhao, L. \& Post, F.\ H. (2007).
Topology-based flow visualization, the state of the art. In \emph{Topology-based Methods in
Visualization}, Springer, Berlin, pp.\ 1--19.
\url{https://doi.org/10.1007/978-3-540-70823-0_1}

\bibitem{lee1991} Lee, S., Lele, S.\ K. \& Moin, P. (1991). Eddy shocklets in decaying
compressible turbulence. \emph{Physics of Fluids A} 3(4), 657--664.
\url{https://doi.org/10.1063/1.858071}

\bibitem{lesnickwright2015} Lesnick, M. \& Wright, M. (2015). Interactive visualization of 2-D
persistence modules. \url{https://arxiv.org/abs/1512.00180}

\bibitem{lesnickwright2022} Lesnick, M. \& Wright, M. (2022). Computing minimal presentations
and bigraded Betti numbers of 2-parameter persistent homology.
\emph{SIAM Journal on Applied Algebra and Geometry} 6(2), 267--298.
\url{https://arxiv.org/abs/1902.05708}

\bibitem{mallat2012} Mallat, S. (2012). Group invariant scattering.
\emph{Communications on Pure and Applied Mathematics} 65(10), 1331--1398.
\url{https://doi.org/10.1002/cpa.21413}

\bibitem{metivier2016gji} M\'etivier, L., Brossier, R., M\'erigot, Q., Oudet, E. \& Virieux,
J. (2016). Measuring the misfit between seismograms using an optimal transport distance:
Application to full waveform inversion. \emph{Geophysical Journal International} 205(1),
345--377. \url{https://doi.org/10.1093/gji/ggw014}

\bibitem{metivier2018geophysics} M\'etivier, L., Allain, A., Brossier, R., M\'erigot, Q.,
Oudet, E. \& Virieux, J. (2018). Optimal transport for mitigating cycle skipping in
full-waveform inversion: A graph-space transform approach. \emph{Geophysics} 83(5),
R515--R540. \url{https://doi.org/10.1190/geo2017-0807.1}

\bibitem{metivier2018eage} M\'etivier, L., Allain, A., Brossier, R., M\'erigot, Q., Oudet, E.
\& Virieux, J. (2018). A graph-space optimal transport approach for full waveform inversion.
\emph{80th EAGE Conference and Exhibition}, 1--5.
\url{https://doi.org/10.3997/2214-4609.201801031}

\bibitem{metivier2019ip} M\'etivier, L., Brossier, R., M\'erigot, Q. \& Oudet, E. (2019). A
graph space optimal transport distance as a generalization of $L^p$ distances: application to
a seismic imaging inverse problem. \emph{Inverse Problems} 35(8), 085001.
\url{https://doi.org/10.1088/1361-6420/ab206f}

\bibitem{muszynski2019} Muszynski, G., Kashinath, K., Kurlin, V., Wehner, M. \& Prabhat
(2019). Topological data analysis and machine learning for recognizing atmospheric river
patterns in large climate datasets. \emph{Geoscientific Model Development} 12(2), 613--628.
\url{https://doi.org/10.5194/gmd-12-613-2019}

\bibitem{peyre2018} Peyre, R. (2018). Comparison between $W_2$ distance and $\dot H^{-1}$
norm, and localization of Wasserstein distance.
\emph{ESAIM: Control, Optimisation and Calculus of Variations} 24(4), 1489--1501.
\url{https://doi.org/10.1051/cocv/2017050}

\bibitem{regaldo2020} Regaldo-Saint Blancard, B., Levrier, F., Allys, E., Bellomi, E. \&
Boulanger, F. (2020). Statistical description of dust polarized emission from the diffuse
interstellar medium: A RWST approach. \emph{Astronomy \& Astrophysics} 642, A217.
\url{https://doi.org/10.1051/0004-6361/202038044}

\bibitem{sadarjoen1999} Sadarjoen, I.\ A. \& Post, F.\ H. (1999). Geometric methods for vortex
detection. In \emph{Data Visualization '99} (VisSym '99), Springer, pp.\ 53--62.

\bibitem{sadarjoen2000} Sadarjoen, I.\ A. \& Post, F.\ H. (2000). Detection, quantification,
and tracking of vortices using streamline geometry. \emph{Computers \& Graphics} 24(3),
333--341. \url{https://doi.org/10.1016/S0097-8493(00)00029-7}

\bibitem{saydjari2021} Saydjari, A.\ K., Portillo, S.\ K.\ N., Slepian, Z., Kahraman, S.,
Burkhart, B. \& Finkbeiner, D.\ P. (2021). Classification of magnetohydrodynamic simulations
using wavelet scattering transforms. \emph{The Astrophysical Journal} 910(2), 122.
\url{https://doi.org/10.3847/1538-4357/abe46d}

\bibitem{schneidervasilyev2010} Schneider, K. \& Vasilyev, O.\ V. (2010). Wavelet methods in
computational fluid dynamics. \emph{Annual Review of Fluid Mechanics} 42, 473--503.
\url{https://doi.org/10.1146/annurev-fluid-121108-145637}

\bibitem{tadmor1987} Tadmor, E. (1987). The numerical viscosity of entropy stable schemes for
systems of conservation laws. I. \emph{Mathematics of Computation} 49(179), 91--103.
\url{https://doi.org/10.1090/S0025-5718-1987-0890255-3}

\bibitem{tadmor2003} Tadmor, E. (2003). Entropy stability theory for difference approximations
of nonlinear conservation laws and related time-dependent problems.
\emph{Acta Numerica} 12, 451--512. \url{https://doi.org/10.1017/S0962492902000156}

\bibitem{taha2015} Taha, A.\ A. \& Hanbury, A. (2015). Metrics for evaluating 3D medical image
segmentation: analysis, selection, and tool. \emph{BMC Medical Imaging} 15, 29.
\url{https://doi.org/10.1186/s12880-015-0068-x}

\bibitem{tierny2018} Tierny, J., Favelier, G., Levine, J.\ A., Gueunet, C. \& Michaux, M.
(2018). The Topology ToolKit. \emph{IEEE Transactions on Visualization and Computer Graphics}
24(1), 832--842. \url{https://doi.org/10.1109/TVCG.2017.2743938}

\bibitem{wernli2008} Wernli, H., Paulat, M., Hagen, M. \& Frei, C. (2008). SAL---A novel
quality measure for the verification of quantitative precipitation forecasts.
\emph{Monthly Weather Review} 136(11), 4470--4487.
\url{https://doi.org/10.1175/2008MWR2415.1}

\bibitem{yang2016} Yang, Y., Engquist, B., Sun, J. \& Froese, B.\ D. (2016). Application of
optimal transport and the quadratic Wasserstein metric to full-waveform inversion.
\url{https://arxiv.org/abs/1612.05075}

\end{thebibliography}

\end{document}